\subsection{One exact finite obstruction map for the whole descent image}

Let $a_g\in\{0,1,2,3,4,5\}$ be the successful exponent
selected by B.2.12 for a normalized target $g\in H$, and
put $l_g=-u_g\bmod3$, where
$\beta_g^{\mathrm{ab}}=(u_g,u_g)$.  Denote the constructed
positive quotient automorphism of B.2.13 by $\alpha_U(g)$.

\paragraph{The full remaining exponent freedom is an actual inner map.}
For any $t\in\widehat{\mathbb Z}$, replace $a_g$ by
$a_g+6t$ while keeping $l_g$ fixed.  The target word $C$
is multiplied on the left by $\tau^{12t}$ and $J_0$
is unchanged.  The two actual source tuples are
\[
\begin{aligned}
 \eta(x^{12t})&=(1,\tau^{12t},\tau^{12t},1,1,1),\\
 \eta(y^{12t})&=(\tau^{12t},1,1,1,1,\tau^{12t}).
\end{aligned}                                                   \tag{B.S24}
\]
Consequently the new parameter tuple is
$y^{12t}w_gx^{-12t}$ in $J$, and the new quotient action
is literally
\[
 \alpha_U(g;a_g+6t)=
               \operatorname{Inn}(x^{12t})\alpha_U(g;a_g).
                                                                  \tag{B.S25}
\]
This inner map has the actual lift
$\operatorname{Inn}(x^{12t})$ in $U^{\mathrm c}$ and in
$B$, and fixes the named $\tau$ exactly.  It therefore
leaves the cusp defect unchanged when the frames are
transported by $x^{12t}$.  For any other choices of frames
its defect class modulo the two images in (B.S11) is
unchanged.  At $g=1$ the corresponding raw full source
word is $y^{12t}x^{-12t}$; its normalization (B.M13) is
the identity.  Thus this freedom has been identified
by an actual source inner morphism.

Define the finite abelian group
\[
 S=\operatorname{im}\kappa_s\,
                    \operatorname{im}\kappa_r\subseteq M,
                 \qquad \overline M=M/S.                      \tag{B.S26}
\]
The additional cusp-centralizer map has trivial image
by B.2.18, so it makes no further quotient in (B.S26).

\begin{result}{B.2.19}{Proved}{The complete finite cusp obstruction for descent}
There is a continuous action $\vartheta:H\to
\operatorname{Aut}(\overline M)$ and a continuous crossed
homomorphism
\[
 \delta:H\longrightarrow\overline M,
       \qquad \delta(gh)=\delta(g)\,
                         \vartheta(g)(\delta(h)),              \tag{B.S27}
\]
defined by the actual cusp defects of positive central
lifts of $\alpha_U(g)$.  Its exact kernel is
\[
                         \delta^{-1}(1)=D_P(H).                 \tag{B.S28}
\]
For the affine action $g\cdot z=\delta(g)\vartheta(g)(z)$
of $H$ on the finite set $\overline M$, one consequently
has
\[
 [H:D_P(H)]=|H\cdot1|=|\delta(H)|\leq|\overline M|.            \tag{B.S29}
\]
The middle expression is the cardinality of the image
set of the crossed homomorphism; no triviality of its
coefficient action is asserted.
\end{result}

\begin{proof}
We first retain the precise composition law for the
quotient actions.  The normalized target frames satisfy
\[
 \beta_{gh}=\alpha_g(\beta_h)\beta_g,\qquad
 A_{gh}=\alpha_g(A_h)A_g,
\]
by the actual automorphism composition and uniqueness
of pure torsion frames in $B$.  In particular $u_g$
is additive and $l_g$ is a homomorphism to $\mathbb Z/3$.
Write $C_g=x^{-a_g}\beta_g$,
$J_g=\beta_g^{-1}A_g r^{l_g}$ and
$\psi_g=\operatorname{Inn}(\beta_g^{-1})\alpha_g$.
The composition of the two quotient source actions has
parameter $w_g\alpha_U(g)(w_h)$.  Its first and fourth
coordinates, using (B.R19), are
\[
 \begin{aligned}
 C_g\psi_g(C_h)&=x^{-(a_g+a_h)}\beta_{gh},\\
 \psi_g(J_h)J_g&=\beta_{gh}^{-1}A_{gh}r^{l_g+l_h}.
 \end{aligned}                                                  \tag{B.S30}
\]
For the first identity use
$\psi_g(x)=C_g^{-1}xC_g$.  For the second use
$\alpha_g(r)=A_g r A_g^{-1}$ and the displayed frame
composition law.  The full six-tuple is determined
from these two coordinates by the source cycles, as
in (B.R12).  It belongs to $J'$, since $w_g,w_h\in J'$
and $\alpha_U(g)$ preserves that subgroup.  The uniqueness
of the successful exponent in B.2.12 therefore gives
$a_{gh}\equiv a_g+a_h\pmod6$.  Thus $a\bmod6$ is a
continuous homomorphism.  Put
$t(g,h)=(a_g+a_h-a_{gh})/6$.  Equations (B.S25) and
(B.S30) give the exact composition comparison
\[
 \alpha_U(g)\alpha_U(h)=
      \operatorname{Inn}(x^{12t(g,h)})\alpha_U(gh).              \tag{B.S31}
\]

For each $g$, choose any positive pure frames of
$\alpha_U(g)$ and form the automatic central
automorphism (B.S9).  Its action on $M$ is independent
of these choices.  Indeed the quotient of two
automorphisms inducing the same action on $U$ changes
each generator of $U^{\mathrm c}$ by a central element
of $M$.  It fixes every commutator and hence the
closed derived subgroup pointwise.  This fixes $M$,
which lies in that derived subgroup.  The inner
comparison in (B.S31) also acts trivially on $M$.
The resulting actions consequently define a genuine
homomorphism $H\to\operatorname{Aut}(M)$.

It preserves $S$.  For example, a positive central
lift sends the commutator defining $\kappa_s(c)$ to
the corresponding commutator with its conjugate of
$s$.  Conjugating back by its pure $s$-frame identifies
this with $\kappa_s(c')$ for an actual
$c'\in C_J(s)$.  Centrality of $M$ removes the inner
conjugation on its value.  The inverse automorphism
gives equality of images; the argument for $r$ is
identical.  The action therefore descends to the
asserted $\vartheta$ on $\overline M$.

By (B.S10)--(B.S11), the defect class of a positive
central lift of $\alpha_U(g)$ is independent of the
chosen pure frames after passing to $\overline M$.
Define $\delta(g)$ to be this class.  It is also
independent of the remaining exponent choice by
(B.S25).  The actions and this defect map are
continuous.  One direct verification uses the compact
space of all triples $(g,\beta_U,A_U)$ satisfying the
positive frame equations for the continuous generator
images of $\alpha_U(g)$.  Its fibres are nonempty by
B.2.13.  Conjugation by a lift of a pure $U$-frame is
a continuous function of that frame, by centrality
of $M$, as proved in B.2.16.  The resulting maps to
the finite action group and to $\overline M$ are
constant on each fibre.  Their compact closed graphs
therefore descend to continuous maps on $H$.

Take positive central lifts for $g$ and $h$.
Their composite sends $\tau$ to
$\tau d_g\widetilde\alpha_g(d_h)$.
Conjugating this composite by the actual
$x^{-12t(g,h)}$ gives a positive central lift above
$\alpha_U(gh)$, by (B.S31).  That conjugation fixes
$\tau$ and $M$, so it changes no cusp defect.
Passing to $\overline M$ proves the crossed law
(B.S27).

Now take an actual normalized preimage $h$ with
$D_P(h)=g$.  Its source automorphism preserves $K$:
it preserves the literal pants kernel, and
conjugation by an automorphism preserves the $B$-core
of that kernel.  Its reduction to $U$ has the two
evaluation words in (B.Q11), with an actual
$a\in\widehat{\mathbb Z}$ and $l=l_g$.
Its parameter lies in $J'$, so B.2.12 forces
$a\equiv a_g\pmod6$.  Equations (B.S24)--(B.S25)
identify its reduced action with
$\operatorname{Inn}(x^{12t})\alpha_U(g)$ for the
actual $t=(a-a_g)/6$.  The full source action induces
a positive $\tau$-fixed automorphism of
$U^{\mathrm c}$.  Removing this literal inner power
of $x$ produces such a lift of $\alpha_U(g)$.
Its defect is one, proving $\delta(g)=1$.

Conversely, $\delta(g)=1$ means that the defect
$d_g$ of a positive central lift lies in $S$.
Choose actual centralizer elements in (B.S11) whose
two images have product $d_g^{-1}$.  Changing the
two $U$-frames by these elements changes the defect
to one and leaves the quotient action $\alpha_U(g)$
unchanged.  The resulting positive central lift
fixes $\tau$ exactly.  B.2.14 supplies a full positive
$\tau$-fixed source automorphism; its exact
commutator normalization changes the target only
by $\operatorname{Inn}(s^b)$, as in B.2.16.
It therefore gives an actual element of $H$ with
descendant $g$.  This proves (B.S28).

The crossed law makes the displayed affine rule an
action on $\overline M$.  Its stabilizer at one is
exactly $\delta^{-1}(1)$.  The orbit--stabilizer
bijection of cosets gives (B.S29).
\end{proof}

The finite module $\overline M$, its coefficient action,
and the image of the crossed homomorphism have not yet
been computed in full.  Equations (B.S27)--(B.S29)
give an exact finite obstruction, rather than a claim
that its image set is trivial or a character with
unproved trivial coefficient action.
